\newcommand{\ANG}[1]{
  {\color{orange} ANG: #1
  }}
\newcommand{\elm}[1]{\textcolor{red}{[EMA: #1]}}
\newcommand{\cmt}[1]{\textcolor{blue}{[CHAR: #1]}}

\newcommand{\SAM}[1]{
  {\color{green} SAM: #1
}}

\newcommand{\EMA}[1]{
  {\color{red} EMA: #1
}}

\definecolor{myBlue}{HTML}{3A7CA5}

% \renewcommand{\ANG}[1]{}
% \renewcommand{\elm}[1]{}
% \renewcommand{\cmt}[1]{}
% \renewcommand{\SAM}[1]{}

\begin{abstract}

    Current economic forces are pushing AI towards higher autonomy. This autonomy is often opposed by greater needs for safety from the models. However, there is no clear agreed understanding of what safe ai constitutes. In this work, we propose a characterization of harmfulness from the point of view of Normative systems. We take ideas from moral philosophy, legal studies along the present works in AI safety, to propose an operationalized notion of safety for both AI to AI and AI to human interactions. In addition, we provide the following contributions: (1) a defeasible deontic logic engine for automatic legal reasoning, aiding the computation of obligations,permissions and prohibitions given certain facts (2) an alternative norms system to constitutional AI that adapts by design to new safety scenarios. In conclusion, we observe our method improving over the classical constitutional AI by \% and better generalization to newer domains by \%.
\end{abstract}

\section{Introduction}

\section{Related Work}

\section{Method}

\section{Limitations}

\onecolumn

\bibliography{aaai2027}
\appendix
\section{Test Appendix}